\changecolours{ScoutTeal}
\section{Contexto Histórico \& Fundamentos}

% ------------------------------------------------------------------------------
\twocolframe[title={A Era dos Sistemas de Arquivos Tradicionais},leftcol=7cm,rightcol=7cm]{%
  \alert{O Cenário das Décadas de 1950 e 1960}\par
  \itemseps[topsepi=2mm,itemsepi=3mm,labelsep=2mm,leftmargini=4mm]
  \begin{itemize}
    \item \textbf{Processamento Baseado em Arquivos:} Antes da criação dos SGBDs, os dados corporativos eram organizados diretamente pelo sistema operacional em arquivos sequenciais ou indexados.
    \item \textbf{Aplicações Monolíticas:} Softwares escritos em linguagens como COBOL, FORTRAN e Assembly precisavam gerenciar o posicionamento físico de ponteiros de leitura e gravação no disco.
    \item \textbf{Fragmentação Setorial:} Cada setor ou departamento da empresa desenvolvia seus próprios programas e formatos de arquivos privativos.
  \end{itemize}
}{%
  \alert{Impacto Operacional}\par
  \itemseps[topsepi=2mm,itemsepi=3mm,labelsep=2mm,leftmargini=4mm]
  \begin{itemize}
    \item \textbf{Alto Custo de Manutenção:} Qualquer alteração na estrutura física de um arquivo exigia recompilar e testar todos os programas que o liam.
    \item \textbf{Dependência Programa-Dado:} A estrutura dos registros (tamanho de campos, tipos) estava intrinsecamente amarrada ao código-fonte.
    \item \textbf{Inviabilidade de Consultas Ad-Hoc:} Responder a uma pergunta gerencial não prevista demandava dias de desenvolvimento de um novo programa.
  \end{itemize}
}

% ------------------------------------------------------------------------------
\twocolframe[title={Limitações dos Arquivos (Parte 1: Estrutura)},leftcol=7cm,rightcol=7cm]{%
  \alert{1. Redundância e Inconsistência}\par
  \itemseps[topsepi=2mm,itemsepi=3mm,labelsep=2mm,leftmargini=4mm]
  \begin{itemize}
    \item As mesmas informações cadastrais eram replicadas em múltiplos arquivos criados por diferentes departamentos.
    \item Atualizações aplicadas em um arquivo não se propagavam automaticamente aos demais, gerando divergência crônica de dados.
    \item Desperdício acentuado de espaço de armazenamento em fita e disco rígido.
  \end{itemize}
}{%
  \alert{2. Acesso Rígido \& 3. Isolamento}\par
  \itemseps[topsepi=2mm,itemsepi=3mm,labelsep=2mm,leftmargini=4mm]
  \begin{itemize}
    \item \textbf{Dificuldade no Acesso:} Ausência de linguagens declarativas padronizadas. Toda nova consulta exigia codificar um novo programa.
    \item \textbf{Isolamento dos Dados:} Arquivos gravados em formatos heterogêneos e dispersos impediam o cruzamento automatizado de relatórios.
  \end{itemize}
}

% ------------------------------------------------------------------------------
\twocolframe[title={Limitações dos Arquivos (Parte 2: Operação)},leftcol=7cm,rightcol=7cm]{%
  \alert{4. Problemas de Integridade}\par
  \itemseps[topsepi=2mm,itemsepi=3mm,labelsep=2mm,leftmargini=4mm]
  \begin{itemize}
    \item Restrições de negócio (ex: número de porta IP válida) ficavam espalhadas pelo código de aplicações individuais.
    \item Caso um novo aplicativo esquecesse de implementar a checagem, corrompia a base inteira compartilhada.
  \end{itemize}
}{%
  \alert{5. Falhas \& 6. Concorrência}\par
  \itemseps[topsepi=2mm,itemsepi=3mm,labelsep=2mm,leftmargini=4mm]
  \begin{itemize}
    \item \textbf{Atomicidade Comprometida:} Quedas de energia durante a gravação deixavam arquivos em estado corrompido pela metade.
    \item \textbf{Acesso Concorrente Descontrolado:} Leituras e gravações simultâneas sem travas causavam perda sistemática de atualizações (\textit{lost updates}).
  \end{itemize}
}

% ------------------------------------------------------------------------------
\twocolframe[title={Definição Formal: Dado, Informação e SGBD},leftcol=7cm,rightcol=7cm]{%
  \alert{Conceitos Canônicos (Date / Elmasri)}\par
  \itemseps[topsepi=2mm,itemsepi=3mm,labelsep=2mm,leftmargini=4mm]
  \begin{itemize}
    \item \textbf{Dado:} Fato bruto conhecido que possui significado implícito e pode ser registrado (ex: uma string ``192.168.1.1'').
    \item \textbf{Informação:} Dado processado, contextualizado e interpretado para tomada de decisão (ex: ``IP do gateway principal'').
    \item \textbf{Metadado:} Descrição formal da estrutura, tipos, tamanhos e restrições dos dados armazenados no banco.
  \end{itemize}
}{%
  \alert{O Que Constitui um SGBD?}\par
  \itemseps[topsepi=2mm,itemsepi=3mm,labelsep=2mm,leftmargini=4mm]
  \begin{itemize}
    \item \textbf{Definição:} Coleção de programas que permite aos usuários criar, manter, consultar e proteger uma base de dados.
    \item \textbf{Catálogo do Sistema:} Repositório central autodescritivo contendo os metadados do esquema.
    \item \textbf{Abstração Operacional:} Isola completamente a lógica da aplicação dos detalhes mecânicos dos discos magnéticos ou SSDs.
  \end{itemize}
}

% ------------------------------------------------------------------------------
\twocolframe[title={Objetivos e Vantagens de um SGBD Moderno},leftcol=7cm,rightcol=7cm]{%
  \alert{Controle \& Centralização}\par
  \itemseps[topsepi=2mm,itemsepi=3mm,labelsep=2mm,leftmargini=4mm]
  \begin{itemize}
    \item \textbf{Redução de Redundância:} Cada fato elementar é registrado em um único local lógico, garantindo consistência.
    \item \textbf{Compartilhamento Seguro:} Múltiplos usuários e processos acessam dados simultâneos com isolamento.
    \item \textbf{Controle de Acesso:} Políticas estritas de autenticação e permissões granulares por usuário.
  \end{itemize}
}{%
  \alert{Confiabilidade \& Flexibilidade}\par
  \itemseps[topsepi=2mm,itemsepi=3mm,labelsep=2mm,leftmargini=4mm]
  \begin{itemize}
    \item \textbf{Garantia Transacional ACID:} Atomicidade, consistência, isolamento e durabilidade em todas as gravações.
    \item \textbf{Consultas Declarativas (SQL):} Otimizador interno deduz o caminho mais veloz para responder cada query.
    \item \textbf{Recuperação Automática:} Mecanismos de diário de bordo (\textit{Write-Ahead Log}) restauram o banco pós-falha.
  \end{itemize}
}
